\scanpage{62}
% archive foliation: 61
\uline{Remarque} Supposons le \unsure{pt préc.} Alors
la catégorie $\unsure{\mathcal P|P}\simeq\unsure{\mathcal P|P}_\bullet$ \deleted{\unsure{est ab.}} admet des $\Sigma$-
noyaux (et des \unsure{Conoyaux}, mais un
bimorphisme n'est pas néc. un iso). \deleted{Je}
\deleted{\unsure{On fera attention que les mono}}
c'est \unsure{tjs|donc} vrai que si $N$ est le noyau
de $A\unsure{-|\longrightarrow}B$, alors $N\unsure{-|\longrightarrow}A$ est un $\Sigma$-
mono, i.e. s'insère dans une suite $\Sigma$-exacte
$0\longrightarrow N\unsure{-|\longrightarrow}A\unsure{-|\longrightarrow}A'\unsure{-|\longrightarrow}0$; mais le mono
$A'\unsure{-|\longrightarrow}B$ ne sera pas en général un
$\Sigma$-mono. Aussi on ne peut interpréter
$(\unsure{\mathcal P|P},\deleted{\unsure{\Psi^0}},C^\bullet,\deleted{\unsure{\Phi^\bullet,\varepsilon}}\inserted{\Sigma})$
comme \unsure{sol. d'un} \unsure{Pb} universel en appliquant tels quels les
développements précédents.

\deleted{Mais, \deleted{\unsure{d'autre part}}\inserted{on peut} généraliser
de façon naturelle, pour une
$B$ $k$-linéaire, la notion d'objet plat,}
\deleted{\unsure{et alors}} le syst. ci-dessus est universel
pour les $B$ \inserted{$k$-additives}, \unsure{munis} d'un complexe
de cochaînes \deleted{$\unsure{C^\bullet}$, et d'une \deleted{\unsure{ss}}\inserted{classe} $\Sigma$
de suites exactes, telle que tt hom
\inserted{$u:A\unsure{-|\longrightarrow}B$} de $B$ ait un noyau tel que
$\operatorname{Ker}u\longrightarrow A$ soit $\in\Sigma$}
$C^\bullet$ et d'une classe
$\Sigma_C$ de suites exactes, telle que tt
morphisme $u:A\unsure{-|\longrightarrow}B$ de $B$ admette
un noyau $K$ tel que $K\unsure{-|\longrightarrow}A$ soit un
$\Sigma_C$-mono : (on \unsure{prend} comme
\unsure{morphismes entre tels syst. ceux qui}
\unsure{transforme|transforment} $C^\bullet$ en $C'^\bullet$ et qui
\unsure{soit|sont} $\Sigma$-exact \uuline{à g}).

